\chapter*{Introduction générale}
\addcontentsline{toc}{chapter}{Introduction générale}
\markboth{Introduction générale}{Introduction générale}

Les systèmes logiciels constituent un patrimoine stratégique pour les organisations. Ils supportent les processus métiers, les échanges avec les clients, la gestion des données et de nombreuses activités critiques. Toutefois, les langages, frameworks, bibliothèques et outils de construction évoluent continuellement. Les applications reposant sur des versions anciennes deviennent alors plus difficiles à maintenir, sécuriser, tester et déployer.

Une migration logicielle ne se limite pas à modifier un numéro de version. Elle exige d'identifier le profil technique de l'application, de définir une trajectoire compatible, d'adapter les dépendances, la configuration et le code source, puis d'exécuter les installations, les compilations et les tests. Chaque échec ouvre un nouveau cycle de diagnostic, de correction et de validation. Réalisée manuellement, cette démarche mobilise des compétences spécialisées, nécessite un temps important et produit des résultats difficiles à reproduire à grande échelle.

Les outils déterministes automatisent efficacement les transformations connues et les contrôles reproductibles. Ils restent toutefois limités lorsque le problème dépend fortement du contexte d'une application. Les grands modèles de langage peuvent, à l'inverse, interpréter des journaux, analyser des extraits de code et proposer des adaptations. Leur fonctionnement probabiliste ne permet cependant pas de leur déléguer sans contrôle les commandes, les modifications et les décisions sensibles.

Le projet de fin d'études réalisé au sein de \textbf{CGI Technologies et Solutions Maroc} répond à cette tension par la conception et le développement d'une \textbf{Migration Factory intelligente, extensible et gouvernée}. La solution combine des services déterministes, des agents spécialisés, un orchestrateur persistant et une supervision humaine. Elle prend en charge la préparation de l'environnement, l'analyse, la planification, la transformation, la validation, la gestion des échecs, la réparation et la production de preuves.

La solution est conçue comme une \textbf{architecture multi-agents hybride et centralement orchestrée}. Les agents possèdent des responsabilités limitées : analyse, review, planification, transformation, réparation et assistance. Ils communiquent au moyen de contrats et d'artefacts. Le backend conserve l'autorité sur l'état, les transitions, les chemins, les commandes et l'application des changements. Le développeur intervient aux points de décision où une approbation, un rejet ou une révision est nécessaire.

La mise en œuvre est déclinée dans deux systèmes et deux dépôts indépendants. Le premier traite les migrations Java/Spring Boot, avec une trajectoire progressive entre Spring Boot 2.1, 2.7, 3.5 et 4.0 ainsi que plusieurs versions de Java. Le second traite Angular 18 à 21 et prépare une extension vers des applications plus anciennes. Les deux systèmes partagent les mêmes principes de gouvernance, d'isolation et de traçabilité, tout en conservant leurs outils, profils, agents et workflows propres.

L'utilisateur principal est le développeur chargé de la migration. L'objectif n'est pas de le remplacer, mais de réduire les opérations répétitives, d'accélérer le diagnostic et de lui fournir un environnement de pilotage capable de présenter l'état, les preuves, les blocages et les actions autorisées.

La problématique générale peut être formulée comme suit :

\begin{quote}
\itshape
Comment concevoir et mettre en œuvre une architecture multi-agents hybride capable de coordonner l'analyse, la planification, la transformation, la validation et la réparation de migrations Java/Spring Boot et Angular, tout en préservant l'intégrité du code source, la traçabilité des décisions, la reproductibilité des opérations et le contrôle du développeur ?
\end{quote}

Pour répondre à cette problématique, le projet poursuit plusieurs objectifs : identifier les agents nécessaires, définir leurs rôles et limites d'autorité, concevoir l'orchestrateur, formaliser les contrats d'échange, intégrer les services déterministes propres à chaque écosystème, mettre en œuvre une réparation gouvernée et vérifier les résultats sur des scénarios représentatifs.

Le rapport est organisé en quatre chapitres :

\begin{itemize}
    \item le \textbf{Chapitre I} présente CGI, le projet, les utilisateurs, les acteurs, les problèmes observés, les besoins fonctionnels et non fonctionnels, les objectifs et la méthodologie ;
    \item le \textbf{Chapitre II} expose les notions d'agent intelligent, de système multi-agents, d'orchestration et de migration, puis compare les principaux travaux et plateformes proches ;
    \item le \textbf{Chapitre III} décrit la conception de l'architecture SMA, les agents, les mécanismes de coordination, l'orchestrateur, les états, les workflows et la supervision humaine ;
    \item le \textbf{Chapitre IV} présente l'architecture technique et l'implémentation effective des deux systèmes, ainsi que les agents, le backend, le cockpit, la persistance, la sécurité et les principales difficultés rencontrées.
\end{itemize}

Une conclusion générale synthétisera les apports, les limites et les perspectives d'industrialisation de la Migration Factory.
